\section{The mesh and the cathode of chamber A}\label{app:gapmap}
\sectionstate{partial}{new 2026-11-23: \texttt{rerun/paper1\_gap/} (\texttt{cellobs.py}, \texttt{joint.py},
\texttt{track\_test.py}, \texttt{track\_grid.py}; \texttt{report.html}); pass 1, the two golden runs of chamber~A.}
\intent{The bow of the board (\cref{sec:bow,app:bow}) and the variation of the gap (\cref{sec:gap}) are two views of
the same two surfaces. One model of the mesh and the cathode describes every observable, the dish of the board is
measured a second time without the reference positions, and a finer map of the drift direction improves the angle
resolution by 16 to 21\,\%.}

\Cref{sec:gap} showed that the drift gap of chamber~A is not uniform, and \cref{sec:bow} showed that its readout board
is dished. These two results are related, because the gap is the distance between two surfaces: the mesh, which sits
on the readout board, and the cathode. In this appendix we describe both surfaces with one model and fit the model to
every observable we have. The fit answers two questions. The first is whether the measurements of the bow and the
measurements of the gap are consistent with each other. The second is whether a better description of the drift
field improves the reconstructed angles.

\subsection{One model for both surfaces}
We describe the height of the mesh by $h(x,y)$ and the height of the cathode by $c(x,y)$, both measured upward along
the reference $z$ axis. Each surface is a quadratic function of the position: a constant, two slopes and three
curvature terms. The gap is the difference of the two surfaces, $g=c-h$.

Every observable of \cref{sec:gap,app:bow} follows from these two surfaces. The position route of \cref{app:bow}
measures the mesh height $h$. The geometric gap and the product of the angle-scale velocity and the drift time measure
the gap $g$. The drift time and the angle-scale velocity also depend on the gap separately, through the local drift
field $V/g$: the drift time is proportional to $g^{1+\beta}$ and the velocity to $g^{-\beta}$, where
$\beta=\mathrm{d}\ln\vd/\mathrm{d}\ln E$.

The drift direction gives a third kind of observable. A drift line leaves the mesh normal to the mesh and ends normal
to the cathode. Averaged over the drift column, its tilt is therefore the mean of the two slopes, $(\partial h/\partial
x+\partial c/\partial x)/2$ in the X view. This tilt shifts the arriving charge sideways in proportion to its drift
time, so it adds an offset to the fitted transverse rate, $w_0=\vd\,(\partial h/\partial x+\partial c/\partial x)/2$.
The same tilt can be measured without any timing. It makes the geometric gap of \cref{sec:gap} read long for tracks of
one sign of $\tan\theta$ and short for the other, by a relative amount equal to the tilt divided by $\tan\theta$. In
the notation of \cref{app:bow}, the ratio $r_f$ is $(1+c''/h'')/2$: it is $1/2$ for a flat cathode and 1 for a cathode
that bows with the board.

We measure seven observables in $4\times4$ cells of the interior (cells of about $85\times75$\,mm), separately in
Run~1 and Run~2. They are the mesh height from the positions (X and Y strips), the transverse-rate offset $w_0$ (X and
Y), the tilt from the sign asymmetry of the geometric gap (X and Y), the geometric gap (X and Y strips), the drift time,
and the angle-scale velocity (X and Y). The fit has 13 shared parameters: the two quadratic surfaces, the mean gap and
$\beta$. The remaining parameters are normalizations, such as the drift velocity of each run and the offset of each
view, which carry no information on the shape. A quadratic surface does not describe the chamber exactly. For this
reason we add to each observable an error in quadrature, chosen so that each class of observables has a reduced
$\chi^2$ of 1. This added error is the size of what the quadratic model does not describe: 0.18\,mm in the mesh height,
0.18\,\textmu m/ns in $w_0$, 0.1\,mm in the geometric gap, 1.9\,\% in the drift time and 1.4\,\% in the velocity.

\subsection{The bow and the gap agree}
The test of consistency is to fit the model to subsets of the observables and to compare the surfaces they give.
\Cref{tab:surfaces} shows the result.

\begin{table}[tbp]
  \centering
  \caption{The mesh and the cathode of chamber~A from subsets of the observables. Curvatures are given as the
  depth of the surface at 150\,mm from the center relative to the center (mm; positive means the center is lower).
  The gap tilt is the slope of the gap along $x$ (mm per 100\,mm). Both runs are fitted together. Subsets: positions
  and the drift tilt ($w_0$ and the sign asymmetry); positions and the gap (geometric gap, drift time, velocity); no
  timing (positions, geometric gap, sign asymmetry); no positions (everything else).}
  \label{tab:surfaces}
  \small
  \begin{tabular}{lccccc}
    \hline
    observables & mesh, $x$ & mesh, $y$ & cathode, $x$ & cathode, $y$ & gap tilt, $x$ \\
    \hline
    all                                     & $0.87\pm0.08$ & $1.48\pm0.10$ & $1.19\pm0.11$ & $0.25\pm0.15$ & $0.39\pm0.08$ \\
    positions, drift tilt & $0.76\pm0.10$ & $1.44\pm0.13$ & $1.52\pm0.23$ & $0.39\pm0.26$ & $0.44\pm0.33$ \\
    positions, gap & $0.76\pm0.10$ & $1.44\pm0.13$ & $1.01\pm0.16$ & $0.18\pm0.25$ & $0.39\pm0.08$ \\
    no timing & $0.77\pm0.10$ & $1.39\pm0.12$ & $0.01\pm0.46$ & $-0.22\pm0.32$ & $0.58\pm0.15$ \\
    no positions                            & $1.02\pm0.12$ & $1.55\pm0.15$ & $1.26\pm0.12$ & $0.29\pm0.15$ & $0.39\pm0.24$ \\
    \hline
  \end{tabular}
\end{table}

Three results follow from \cref{tab:surfaces}. First, the dish of the board is measured twice. The position route sees
it through the reference tracks. Without any position information, the drift tilt measures $h+c$ and the gap measures
$c-h$, so together they also determine the mesh. The two measurements agree: the mesh lies 0.8--1.0\,mm (along $x$)
and 1.4--1.6\,mm (along $y$) lower at the center than 150\,mm out. When we let the position route carry a separate dish
of the reference telescope, the fit gives $-0.25\pm0.16$ and $-0.05\pm0.20$\,mm, which is consistent with zero.

Second, along $y$ every subset finds a nearly flat cathode. Its depth at 150\,mm is between $-0.2$ and $0.4$\,mm,
against 1.4--1.6\,mm for the mesh. The gap along $y$ is therefore longest at the center, by about 1.2\,mm relative to a
point 150\,mm out. This is the dome of \cref{sec:gap}, and the drift tilt alone predicts it as well as the gap
measurements do.

Third, along $x$ the cathode follows the board. The gap then has little curvature along $x$, and it increases steadily
with $x$ by about 0.4\,mm per 100\,mm. The only subset that disagrees is the one without timing, which finds a flat
cathode along $x$ ($2.5\sigma$ from the full fit). The geometric gap with the X strips reaches only the central 220\,mm
in $x$, so this subset constrains the curvature along $x$ weakly. In addition, the drift tilt along $x$ rises faster
toward the edges than a quadratic dish allows: its gradient is 0.39\,\textmu m/ns per 100\,mm over the interior and
0.31 in the central 200\,mm. With the full interior, $r_f$ along $x$ is therefore larger than 1 (1.2). We do not
interpret this as a cathode that bows more than the board. It indicates that the board is not a pure paraboloid, which
\cref{app:bow} also notes.

The fit gives $\beta=1.0\pm0.2$. Magboltz gives 0.68 at this field and water content, and the drift-time scan alone
gives 0.87 (\cref{sec:vE}). Fixing $\beta$ at 0.68 raises the $\chi^2$ by 5, and fixing it at 0.87 raises it by 0.7. The
shape of the surfaces does not depend on this choice (the cathode depths change by less than 0.2\,mm).

\begin{figure}[tbp]
  \centering
  \includegraphics[width=\linewidth]{gap_surfaces}
  \caption{The two surfaces of chamber~A and the angle resolution. (a) Drift gap as a function of $x$ at mid-gap,
  averaged over $y$, and (b) as a function of $y$, averaged over $x$, both runs combined: geometric gap with the X
  strips (filled circles) and the Y strips (open circles; relative, placed on 28.19\,mm), and the product of the
  angle-scale velocity and the drift time (squares). Line: the fit of the mesh and cathode model to all observables.
  (c) Angle resolution (s68 of $\theta-\theta_\mathrm{ref}$) with the pass-1 conversion (left) and with a map of the
  drift-tilt offset and the velocity made on the other run, as a function of the cell size of the map. X solid, Y
  dashed; Run~1 filled, Run~2 open.}
  \label{fig:gapsurf}
\end{figure}

\subsection{A finer map of the drift direction improves the angles}
The reconstruction converts the fitted transverse rate into an angle with the bow block of \cref{sec:bow}. This block
corrects the drift tilt with a quadratic surface and the angle scale with a linear term. The added errors above show
that the drift tilt has structure that a quadratic surface does not describe, with an amplitude of 0.18\,\textmu m/ns.
This corresponds to about 0.005 in $\tan\theta$, or 0.3\textdegree.

To test whether this structure can be corrected, we make a map of the transverse-rate offset and of the velocity on
one run and apply it to the other run. The map is a grid of cells over the interior, interpolated between the cell
centers, and it is indexed by the reconstructed position of the track, not by the reference. Only the two constants
of the conversion (the mean offset and the mean velocity) are measured again on the run itself. Since the map and the
test use different runs, the gain cannot come from fitting the noise of the test sample.

\Cref{fig:gapsurf}c and \cref{tab:anglemap} show the result. The smooth two-surface model gives a small gain on the
X strips (0.03--0.04\textdegree). A map with cells of about 40\,mm improves the angle resolution by 0.15--0.17\textdegree,
which is 16 to 21\,\% on both views and in both runs. Most of the gain comes from the offset map. The velocity map
adds a further 0.02--0.04\textdegree. With cells smaller than about 40\,mm, the map made on Run~1 becomes noisy and the gain
on Run~2 falls.

\begin{table}[tbp]
  \centering
  \caption{Angle resolution of chamber~A (s68 of $\theta-\theta_\mathrm{ref}$, degrees), single tracks in the
  interior with $|\tan\theta_\mathrm{ref}|\le0.6$. Each map is made on the other run.}
  \label{tab:anglemap}
  \begin{tabular}{lcccc}
    \hline
    angle conversion & X, Run~1 & X, Run~2 & Y, Run~1 & Y, Run~2 \\
    \hline
    pass 1 (quadratic bow block)            & 0.905 & 0.929 & 0.796 & 0.809 \\
    two-surface model (quadratic)           & 0.878 & 0.888 & 0.797 & 0.791 \\
    map, 85\,mm cells                       & 0.816 & 0.837 & 0.717 & 0.730 \\
    map, 57\,mm cells                       & 0.764 & 0.790 & 0.661 & 0.679 \\
    map, 42\,mm cells                       & 0.736 & 0.776 & 0.627 & 0.651 \\
    map, 42\,mm cells, offset only          & 0.779 & 0.807 & 0.653 & 0.675 \\
    \hline
  \end{tabular}
\end{table}

The structure the map corrects is two-dimensional and smooth. It is not a timing offset between front-end chips. A
chip reads 64 strips (50\,mm), so a chip offset would appear as steps in the view's own coordinate. The offset varies
smoothly along that coordinate, without steps at the chip boundaries. A map in that coordinate alone recovers between
a fifth and a third of the gain (0.852 and 0.873\textdegree\ on X, 0.762 and 0.767\textdegree\ on Y). A map in the other coordinate
alone recovers none of it. We therefore attribute the structure to the shape of the drift field, which follows the
shape of the board and of the cathode on scales of about 50\,mm.

The map is a property of one chamber in one mounting, and it is measured with the reference tracks. At n\_TOF the
boards stand vertical, so the part of the bow caused by gravity is absent, and we have no reference tracks there. We
have not tested whether any part of the bench map applies at n\_TOF.
\askDylan{Ship a per-chamber 2-D map of $w_0$ and $\vd$ (42\,mm cells, the bundle's \texttt{v\_map} mechanism of the FA
spec) in a bench re-finalization? It needs a cross-run map per chamber and mounting; chambers with one good run would
need a split-in-time version of the test. Numbers: \texttt{rerun/paper1\_gap/results/track\_grid.json}.}
